%% ==================================================================
%% TOAN12_C2_BAI7.tex -- Toán 12, Chương II, Bài 7
%% Cấu trúc: Lí thuyết trọng tâm -> Ví dụ minh hoạ (chừa chỗ trống) -> Bài tập luyện tập
%% ==================================================================
\baihoc{7}{Hệ trục toạ độ trong không gian}

\lithuyet

\begin{dinhnghia}[Hệ trục toạ độ trong không gian]
\begin{minipage}[c]{.6\linewidth}
Trong không gian, ba trục $Ox$, $Oy$, $Oz$ đôi một vuông góc với nhau tại gốc $O$ của mỗi trục. Gọi $\vec i$, $\vec j$, $\vec k$ lần lượt là các vectơ đơn vị trên các trục $Ox$, $Oy$, $Oz$.
\begin{itemize}
  \item Hệ ba trục như vậy được gọi là \textbf{hệ trục toạ độ Descartes vuông góc} $Oxyz$, hay đơn giản là \textbf{hệ toạ độ} $Oxyz$.
  \item Điểm $O$ được gọi là \emph{gốc toạ độ}.
  \item Các mặt phẳng $(Oxy)$, $(Oyz)$, $(Ozx)$ đôi một vuông góc với nhau được gọi là các \emph{mặt phẳng toạ độ}.
\end{itemize}
Không gian với hệ toạ độ $Oxyz$ còn được gọi là \emph{không gian} $Oxyz$.
\end{minipage}\hfill
\begin{minipage}[c]{.37\linewidth}\centering
\begin{tikzpicture}[scale=1.15,font=\small,x={(-0.55cm,-0.45cm)},y={(1cm,0cm)},z={(0cm,1cm)}]
  \fill[dcnen,opacity=.9] (0,0,0)--(2,0,0)--(2,2,0)--(0,2,0)--cycle;
  \fill[dcvang,opacity=.8] (0,0,0)--(0,2,0)--(0,2,1.8)--(0,0,1.8)--cycle;
  \fill[dcnhan!12,opacity=.8] (0,0,0)--(2,0,0)--(2,0,1.8)--(0,0,1.8)--cycle;
  \draw[truc] (0,0,0)--(2.6,0,0) node[below left]{$x$};
  \draw[truc] (0,0,0)--(0,2.6,0) node[right]{$y$};
  \draw[truc] (0,0,0)--(0,0,2.2) node[above]{$z$};
  \draw[vecto,dcnhan] (0,0,0)--(1,0,0) node[left]{$\vec i$};
  \draw[vecto,dcnhan] (0,0,0)--(0,1,0) node[above]{$\vec j$};
  \draw[vecto,dcnhan] (0,0,0)--(0,0,1) node[left]{$\vec k$};
  \node[above left] at (0,0,0){$O$};
  \node[text=dcchinh,font=\scriptsize] at (1.2,1.4,0){$(Oxy)$};
  \node[text=dcchinh,font=\scriptsize] at (0,1.4,1.3){$(Oyz)$};
\end{tikzpicture}
\end{minipage}
\end{dinhnghia}

\begin{dinhnghia}[Toạ độ của điểm]
\begin{minipage}[c]{.6\linewidth}
Trong không gian $Oxyz$, cho một điểm $M$ tuỳ ý. Bộ ba số $(x;y;z)$ duy nhất sao cho $\vt{OM}=x\vec i+y\vec j+z\vec k$ được gọi là \textbf{toạ độ} của điểm $M$ đối với hệ toạ độ $Oxyz$. Khi đó, ta viết $M=(x;y;z)$ hoặc $M(x;y;z)$, trong đó $x$ là \textbf{hoành độ}, $y$ là \textbf{tung độ} và $z$ là \textbf{cao độ} của $M$.
\end{minipage}\hfill
\begin{minipage}[c]{.37\linewidth}\centering
\begin{tikzpicture}[scale=1,font=\small,x={(-0.55cm,-0.45cm)},y={(1cm,0cm)},z={(0cm,1cm)}]
  \coordinate (O) at (0,0,0); \coordinate (A) at (1.6,0,0); \coordinate (B) at (0,2.2,0); \coordinate (C) at (0,0,1.6);
  \coordinate (D) at (1.6,2.2,0); \coordinate (E) at (0,2.2,1.6); \coordinate (F) at (1.6,0,1.6); \coordinate (M) at (1.6,2.2,1.6);
  \draw[truc] (O)--(2.4,0,0) node[below left]{$x$}; \draw[truc] (O)--(0,3,0) node[right]{$y$}; \draw[truc] (O)--(0,0,2.3) node[above]{$z$};
  \draw[dcdam] (A)--(D)--(B) (A)--(F)--(C)--(E)--(B) (F)--(M)--(E) (D)--(M);
  \draw[vecto,dcnhan] (O)--(M);
  \node[diem] at (M){}; \node[above right] at (M){$M$};
  \node[left] at (A){$A$}; \node[below] at (B){$B$}; \node[left] at (C){$C$};
  \node[below] at (D){$D$}; \node[above] at (E){$E$}; \node[left] at (F){$F$}; \node[above left] at (O){$O$};
\end{tikzpicture}
\end{minipage}
\end{dinhnghia}

\begin{nhanxet}
Nếu điểm $M$ có toạ độ $(x;y;z)$ đối với hệ toạ độ $Oxyz$ thì:
\begin{itemize}
  \item Hình chiếu vuông góc của $M$ trên các trục $Ox$, $Oy$ và $Oz$ có toạ độ lần lượt là $(x;0;0)$, $(0;y;0)$ và $(0;0;z)$.
  \item Hình chiếu vuông góc của $M$ trên các mặt phẳng $(Oxy)$, $(Oyz)$ và $(Ozx)$ có toạ độ lần lượt là $(x;y;0)$, $(0;y;z)$ và $(x;0;z)$.
\end{itemize}
\end{nhanxet}

\begin{dinhnghia}[Toạ độ của vectơ]
Trong không gian $Oxyz$, cho vectơ $\vec a$ tuỳ ý. Bộ ba số $(x;y;z)$ duy nhất sao cho $\vec a=x\vec i+y\vec j+z\vec k$ được gọi là \textbf{toạ độ của vectơ} $\vec a$ đối với hệ toạ độ $Oxyz$. Khi đó, ta viết $\vec a=(x;y;z)$ hoặc $\vec a(x;y;z)$.
\end{dinhnghia}

\begin{nhanxet}
\begin{itemize}
  \item Toạ độ của vectơ $\vec a$ cũng là toạ độ của điểm $M$ sao cho $\vt{OM}=\vec a$.
  \item Trong không gian, cho hai vectơ $\vec a=(x;y;z)$ và $\vec b=(x';y';z')$. Khi đó, $\vec a=\vec b$ nếu và chỉ nếu $x=x'$, $y=y'$, $z=z'$.
  \item $\vec i=(1;0;0)$, $\vec j=(0;1;0)$, $\vec k=(0;0;1)$.
\end{itemize}
\end{nhanxet}

\begin{congthuc}[Toạ độ của vectơ theo toạ độ hai đầu mút]
Trong không gian $Oxyz$, cho hai điểm $M(x_M;y_M;z_M)$ và $N(x_N;y_N;z_N)$. Khi đó:
\[\vt{MN}=(x_N-x_M;\ y_N-y_M;\ z_N-z_M).\]
\end{congthuc}

%% ------------------------------------------------ VÍ DỤ
\viduminhhoa
\subsubsection{Xác định toạ độ điểm, toạ độ vectơ}
\begin{vidu}
Cho hình hộp chữ nhật $OABC.O'A'B'C'$ có $A(4;0;0)$, $C(0;3;0)$, $O'(0;0;2)$. Tìm toạ độ điểm $B'$ và vectơ $\vt{AC'}$.
\begin{loigiai}
\giai $B(4;3;0)$ nên $B'(4;3;2)$; $C'(0;3;2)$ nên $\vt{AC'}=(-4;3;2)$.
\end{loigiai}
\end{vidu}
\begin{vidu}
Cho điểm $M(3;-2;4)$. Tìm toạ độ hình chiếu của $M$ lên $(Oxy)$, lên trục $Oz$, điểm đối xứng với $M$ qua gốc $O$ và qua trục $Ox$.
\begin{loigiai}
\giai Hình chiếu lên $(Oxy)$: $(3;-2;0)$; lên trục $Oz$: $(0;0;4)$. Đối xứng qua $O$: $(-3;2;-4)$; qua trục $Ox$: $(3;2;-4)$.
\end{loigiai}
\end{vidu}
\subsubsection{Toạ độ hoá bài toán thực tiễn}
\begin{vidu}
Một căn phòng dạng hình hộp chữ nhật dài $6$ m, rộng $4$ m, cao $3$ m. Chọn hệ trục $Oxyz$ có gốc $O$ là một góc sàn, các trục lần lượt theo chiều dài, chiều rộng, chiều cao. Một camera đặt ở góc trần đối diện với $O$, một chiếc đèn ở chính giữa trần. Tìm toạ độ camera, đèn và vectơ từ camera đến đèn.
\begin{loigiai}
\giai Camera ở $C(6;4;3)$, đèn ở $E(3;2;3)$; $\vt{CE}=(-3;-2;0)$.
\end{loigiai}
\end{vidu}

\begin{viduds}
Trong không gian $Oxyz$, cho hình hộp chữ nhật $OABC.O'A'B'C'$ có $A(2;0;0)$, $C(0;3;0)$, $O'(0;0;4)$.
\begin{dungsai}
  \y Toạ độ điểm $B$ là $(2;3;0)$. \oDS
  \y Toạ độ điểm $B'$ là $(2;3;4)$. \oDS
  \y $\vt{AC'}=(-2;3;4)$. \oDS
  \y Hình chiếu vuông góc của $B'$ lên mặt phẳng $(Oyz)$ là điểm $A'$. \oDS
\end{dungsai}
\begin{loigiai}
\tuduy \emph{Vẽ nhanh hình hộp và ghi toạ độ cả $8$ đỉnh} trước khi đọc các ý: các đỉnh đáy dưới có $z=0$, đáy trên có $z=4$. Hình chiếu lên $(Oyz)$ thì ``xoá'' hoành độ ($x=0$).
\giai $B(2;3;0)$, $B'(2;3;4)$, $C'(0;3;4)$, $A'(2;0;4)$. a) \textbf{Đúng}. b) \textbf{Đúng}. c) \textbf{Đúng}. d) \textbf{Sai}: hình chiếu của $B'$ lên $(Oyz)$ là $(0;3;4)$, chính là $C'$.
\end{loigiai}
\end{viduds}

\subsubsection{Bài toán thực tế}
\begin{vidu}
Một phòng học có dạng hình hộp chữ nhật với chiều dài $8$ m, chiều rộng $4$ m, chiều cao $4$ m. Chọn hệ trục $Oxyz$ có gốc $O$ tại một góc sàn, các trục $Ox$, $Oy$, $Oz$ lần lượt theo chiều dài, chiều rộng, chiều cao của phòng. Một bóng đèn treo ở chính giữa trần nhà, một camera gắn ở góc trần đối diện với góc $O$ qua tâm phòng. Tìm khoảng cách (m) từ bóng đèn đến camera (làm tròn kết quả đến hàng phần trăm).
\hinh{\begin{tikzpicture}[scale=0.8,font=\footnotesize,line join=round,line cap=round,>=Stealth]
\draw (-1.080,-0.810)--(2.920,-0.810)--(2.920,1.590)--(-1.080,1.590)--cycle;
\draw (-1.080,1.590)--(0.000,2.400)--(4.000,2.400)--(2.920,1.590);
\draw[dashed] (-1.080,-0.810)--(0.000,0.000)--(4.000,0.000) (0.000,0.000)--(0.000,2.400);
\draw (4.000,0.000)--(4.000,2.400);
\draw (2.920,-0.810)--(4.000,0.000);
\node[below left] at (0.000,0.000){$O$};
\draw[->,gray] (-1.080,-0.810)--(-1.368,-1.026) node[below left]{$x$}; \draw[->,gray] (4.000,0.000)--(4.800,0.000) node[right]{$y$}; \draw[->,gray] (0.000,2.400)--(0.000,3.100) node[above]{$z$};
\draw (1.460,1.995)--(1.460,1.695); \shade[ball color=gray!30] (1.460,1.615) circle (.12); \node[right] at (1.460,1.615){đèn};
\filldraw[fill=gray!60] (2.920,1.590) circle (.08) node[right]{camera};

\end{tikzpicture}}

\begin{loigiai}
\giai $(4;2;4)$, $(8;4;4)$; khoảng cách $=2 \sqrt{5}$. Kết quả: $4{,}47$.
\end{loigiai}
\end{vidu}
\begin{vidu}
Ba trạm quan trắc được đặt tại các vị trí $A(5;8;3)$, $B(5;0;1)$, $C(-1;7;2)$ trong hệ trục $Oxyz$ (đơn vị: km). Một trạm điều phối được đặt tại trọng tâm $G$ của tam giác $ABC$. Tìm hoành độ của điểm $G$.
\hinh{\begin{tikzpicture}[scale=0.85,font=\footnotesize,line join=round,line cap=round,>=Stealth]
\fill[gray!10] (0,0)--(4.5,.4)--(2.2,2.6)--cycle; \draw[thick] (0,0) node[below left]{$A$}--(4.5,.4) node[below right]{$B$}--(2.2,2.6) node[above]{$C$}--cycle;
\draw[dashed] (0,0)--(3.35,1.5) (4.5,.4)--(1.1,1.3) (2.2,2.6)--(2.25,.2); \fill (2.233,1) circle (.06) node[right]{$G$};
\foreach \p in {(0,0),(4.5,.4),(2.2,2.6)}{\filldraw[fill=white] \p circle (.07);}
\end{tikzpicture}}

\begin{loigiai}
\giai $x_G=\dfrac{5+5-1}{3}=3$. Kết quả: $3$.
\end{loigiai}
\end{vidu}
\begin{vidu}
Ba trạm quan trắc được đặt tại các vị trí $A(-1;5;1)$, $B(2;-8;2)$, $C(2;-6;0)$ trong hệ trục $Oxyz$ (đơn vị: km). Một trạm điều phối được đặt tại trọng tâm $G$ của tam giác $ABC$. Tìm khoảng cách (km) giữa hai trạm $A$ và $B$ (làm tròn kết quả đến hàng phần mười).
\hinh{\begin{tikzpicture}[scale=0.85,font=\footnotesize,line join=round,line cap=round,>=Stealth]
\fill[gray!10] (0,0)--(4.5,.4)--(2.2,2.6)--cycle; \draw[thick] (0,0) node[below left]{$A$}--(4.5,.4) node[below right]{$B$}--(2.2,2.6) node[above]{$C$}--cycle;
\draw[dashed] (0,0)--(3.35,1.5) (4.5,.4)--(1.1,1.3) (2.2,2.6)--(2.25,.2); \fill (2.233,1) circle (.06) node[right]{$G$};
\foreach \p in {(0,0),(4.5,.4),(2.2,2.6)}{\filldraw[fill=white] \p circle (.07);}
\end{tikzpicture}}

\begin{loigiai}
\giai $AB=\sqrt{179}$. Kết quả: $13{,}4$.
\end{loigiai}
\end{vidu}
\begin{viduds}
Một phòng học có dạng hình hộp chữ nhật với chiều dài $6$ m, chiều rộng $6$ m, chiều cao $3$ m. Chọn hệ trục $Oxyz$ có gốc $O$ tại một góc sàn, các trục $Ox$, $Oy$, $Oz$ lần lượt theo chiều dài, chiều rộng, chiều cao của phòng. Một bóng đèn treo ở chính giữa trần nhà, một camera gắn ở góc trần đối diện với góc $O$ qua tâm phòng. Xét tính đúng sai của các mệnh đề sau:
\hinh{\begin{tikzpicture}[scale=0.8,font=\footnotesize,line join=round,line cap=round,>=Stealth]
\draw (-1.080,-0.810)--(2.920,-0.810)--(2.920,1.590)--(-1.080,1.590)--cycle;
\draw (-1.080,1.590)--(0.000,2.400)--(4.000,2.400)--(2.920,1.590);
\draw[dashed] (-1.080,-0.810)--(0.000,0.000)--(4.000,0.000) (0.000,0.000)--(0.000,2.400);
\draw (4.000,0.000)--(4.000,2.400);
\draw (2.920,-0.810)--(4.000,0.000);
\node[below left] at (0.000,0.000){$O$};
\draw[->,gray] (-1.080,-0.810)--(-1.368,-1.026) node[below left]{$x$}; \draw[->,gray] (4.000,0.000)--(4.800,0.000) node[right]{$y$}; \draw[->,gray] (0.000,2.400)--(0.000,3.100) node[above]{$z$};
\draw (1.460,1.995)--(1.460,1.695); \shade[ball color=gray!30] (1.460,1.615) circle (.12); \node[right] at (1.460,1.615){đèn};
\filldraw[fill=gray!60] (2.920,1.590) circle (.08) node[right]{camera};

\end{tikzpicture}}

\begin{dungsai}
  \y Bóng đèn cách camera $3 \sqrt{2}$ m. \oDS
  \y Tổng các toạ độ của bóng đèn bằng $10$. \oDS
  \y Khoảng cách từ $O$ đến camera bằng $9$ m. \oDS
  \y Góc giữa $OC$ và mặt sàn xấp xỉ $17^\circ$. \oDS
\end{dungsai}
\begin{loigiai}
\giai Camera tại $(6;6;3)$, khoảng cách $=\sqrt{6^2+6^2+3^2}=9$. $\overrightarrow{OC}=(6;6;3)$, $\sin\varphi=\dfrac{3}{9}$, $\varphi\approx 19^\circ$. $(3;3;3)$, $(6;6;3)$; khoảng cách $=3 \sqrt{2}$. Đèn ở chính giữa trần: $(3;3;3)$. a) Đúng. b) Sai: giá trị đúng là $9$. c) Đúng. d) Sai: giá trị đúng là $19$.
\end{loigiai}
\end{viduds}

%% ------------------------------------------------ BÀI TẬP
\luyentap
\phanbt{Trắc nghiệm nhiều phương án lựa chọn}
\begin{bai}Trong không gian $Oxyz$, vectơ $\vec u=3\vec i-2\vec k$ có toạ độ là
\pa{$(3;0;-2)$}{$(3;-2;0)$}{$(0;3;-2)$}{$(3;-2;1)$}\dapso{A.}\end{bai}
\begin{bai}Hình chiếu vuông góc của điểm $M(2;-1;5)$ lên mặt phẳng $(Ozx)$ là
\pa{$(2;0;5)$}{$(0;-1;5)$}{$(2;-1;0)$}{$(0;0;5)$}\dapso{A.}\end{bai}
\begin{bai}Điểm đối xứng với $A(1;-2;3)$ qua mặt phẳng $(Oxy)$ là
\pa{$(1;-2;-3)$}{$(-1;2;3)$}{$(-1;-2;3)$}{$(1;2;3)$}\dapso{A.}\end{bai}
\begin{bai}Trong không gian $Oxyz$, cho $\vt{OM}=2\vec j-3\vec k+\vec i$. Toạ độ điểm $M$ là
\pa{$(1;2;-3)$}{$(2;-3;1)$}{$(2;1;-3)$}{$(-3;2;1)$}\dapso{A. Chú ý thứ tự $\vec i,\vec j,\vec k$.}\end{bai}
\begin{bai}Điểm nào dưới đây thuộc trục $Oy$?
\pa{$(0;-3;0)$}{$(-3;0;0)$}{$(0;0;-3)$}{$(1;-3;0)$}\dapso{A.}\end{bai}
\begin{bai}Cho hình lập phương $ABCD.A'B'C'D'$ cạnh $2$. Chọn hệ trục toạ độ sao cho $A(0;0;0)$, $B(2;0;0)$, $D(0;2;0)$, $A'(0;0;2)$. Toạ độ tâm $I$ của mặt $BCC'B'$ là
\pa{$(2;1;1)$}{$(1;1;2)$}{$(1;2;1)$}{$(1;1;1)$}\dapso{A. $I$ là trung điểm $BC'$ với $C'(2;2;2)$.}\end{bai}

\begin{bai}Trong không gian $Oxyz$ (đơn vị: mét), mặt đất là mặt phẳng $(Oxy)$, hai flycam đang bay ở các vị trí $A(-9;-5;15)$ và $B(2;-3;16)$. Trung điểm $M$ của đoạn $AB$ có cao độ bằng
\hinh{\begin{tikzpicture}[scale=0.85,font=\footnotesize,line join=round,line cap=round,>=Stealth]
\draw[->,gray] (0.000,0.000)--(-1.440,-1.080) node[below left]{$x$};
\draw[->,gray] (0.000,0.000)--(5.000,0.000) node[right]{$y$};
\draw[->,gray] (0.000,0.000)--(0.000,3.300) node[above]{$z$};
\node[below right,gray] at (0.000,0.000){$O$};
\fill[gray!10] (-0.500,0.000)--(-1.760,-0.945)--(3.240,-0.945)--(4.500,0.000)--cycle;
\draw[dashed] (0.300,1.525)--(0.300,-0.675); \fill (0.300,-0.675) circle (.04);
\filldraw[fill=gray!40] (0.300,1.525) ++(-.18,0) rectangle ++(.36,.1); \node[above] at (0.300,1.525){$A$};
\draw[dashed] (3.240,2.530)--(3.240,-0.270); \fill (3.240,-0.270) circle (.04);
\filldraw[fill=gray!40] (3.240,2.530) ++(-.18,0) rectangle ++(.36,.1); \node[above] at (3.240,2.530){$B$};
\draw[thick,dotted] (0.300,1.525)--(3.240,2.530);
\end{tikzpicture}}

\pa{$\dfrac{33}{2}$.}{$\dfrac{31}{2}$.}{$\dfrac{31}{4}$.}{$\dfrac{29}{2}$.}
\dapso{B. $M(-\dfrac{7}{2};-4;\dfrac{31}{2})$.}\end{bai}
\begin{bai}Một phòng học có dạng hình hộp chữ nhật với chiều dài $4$ m, chiều rộng $3$ m, chiều cao $4$ m. Chọn hệ trục $Oxyz$ có gốc $O$ tại một góc sàn, các trục $Ox$, $Oy$, $Oz$ lần lượt theo chiều dài, chiều rộng, chiều cao của phòng. Một bóng đèn treo ở chính giữa trần nhà, một camera gắn ở góc trần đối diện với góc $O$ qua tâm phòng. Góc (độ) giữa đường thẳng nối góc $O$ với camera và mặt sàn, làm tròn đến hàng đơn vị, bằng
\hinh{\begin{tikzpicture}[scale=0.8,font=\footnotesize,line join=round,line cap=round,>=Stealth]
\draw (-1.080,-0.810)--(2.920,-0.810)--(2.920,1.590)--(-1.080,1.590)--cycle;
\draw (-1.080,1.590)--(0.000,2.400)--(4.000,2.400)--(2.920,1.590);
\draw[dashed] (-1.080,-0.810)--(0.000,0.000)--(4.000,0.000) (0.000,0.000)--(0.000,2.400);
\draw (4.000,0.000)--(4.000,2.400);
\draw (2.920,-0.810)--(4.000,0.000);
\node[below left] at (0.000,0.000){$O$};
\draw[->,gray] (-1.080,-0.810)--(-1.368,-1.026) node[below left]{$x$}; \draw[->,gray] (4.000,0.000)--(4.800,0.000) node[right]{$y$}; \draw[->,gray] (0.000,2.400)--(0.000,3.100) node[above]{$z$};
\draw (1.460,1.995)--(1.460,1.695); \shade[ball color=gray!30] (1.460,1.615) circle (.12); \node[right] at (1.460,1.615){đèn};
\filldraw[fill=gray!60] (2.920,1.590) circle (.08) node[right]{camera};

\end{tikzpicture}}

\pa{$39$.}{$29$.}{$78$.}{$34$.}
\dapso{A. $\overrightarrow{OC}=(4;3;4)$, $\sin\varphi=\dfrac{4}{\sqrt{41}}$, $\varphi\approx 39^\circ$.}\end{bai}
\begin{bai}Trong không gian $Oxyz$ (đơn vị: mét), mặt đất là mặt phẳng $(Oxy)$, hai flycam đang bay ở các vị trí $A(6;7;5)$ và $B(1;5;5)$. Khoảng cách từ flycam ở $B$ đến điểm $H$ trên mặt đất nằm ngay dưới $A$ bằng
\hinh{\begin{tikzpicture}[scale=0.85,font=\footnotesize,line join=round,line cap=round,>=Stealth]
\draw[->,gray] (0.000,0.000)--(-1.440,-1.080) node[below left]{$x$};
\draw[->,gray] (0.000,0.000)--(5.000,0.000) node[right]{$y$};
\draw[->,gray] (0.000,0.000)--(0.000,3.300) node[above]{$z$};
\node[below right,gray] at (0.000,0.000){$O$};
\fill[gray!10] (-0.500,0.000)--(-1.760,-0.945)--(3.240,-0.945)--(4.500,0.000)--cycle;
\draw[dashed] (0.300,1.525)--(0.300,-0.675); \fill (0.300,-0.675) circle (.04);
\filldraw[fill=gray!40] (0.300,1.525) ++(-.18,0) rectangle ++(.36,.1); \node[above] at (0.300,1.525){$A$};
\draw[dashed] (3.240,2.530)--(3.240,-0.270); \fill (3.240,-0.270) circle (.04);
\filldraw[fill=gray!40] (3.240,2.530) ++(-.18,0) rectangle ++(.36,.1); \node[above] at (3.240,2.530){$B$};
\draw[thick,dotted] (0.300,1.525)--(3.240,2.530);
\end{tikzpicture}}

\pa{$2 \sqrt{6}\,\text{ m}$.}{$\dfrac{3 \sqrt{6}}{2}\,\text{ m}$.}{$\dfrac{9 \sqrt{6}}{2}\,\text{ m}$.}{$3 \sqrt{6}\,\text{ m}$.}
\dapso{D. $H(6;7;0)$, $BH=3 \sqrt{6}$.}\end{bai}
\begin{bai}Một phòng học có dạng hình hộp chữ nhật với chiều dài $4$ m, chiều rộng $4$ m, chiều cao $3$ m. Chọn hệ trục $Oxyz$ có gốc $O$ tại một góc sàn, các trục $Ox$, $Oy$, $Oz$ lần lượt theo chiều dài, chiều rộng, chiều cao của phòng. Một bóng đèn treo ở chính giữa trần nhà, một camera gắn ở góc trần đối diện với góc $O$ qua tâm phòng. Tổng các toạ độ của bóng đèn bằng
\hinh{\begin{tikzpicture}[scale=0.8,font=\footnotesize,line join=round,line cap=round,>=Stealth]
\draw (-1.080,-0.810)--(2.920,-0.810)--(2.920,1.590)--(-1.080,1.590)--cycle;
\draw (-1.080,1.590)--(0.000,2.400)--(4.000,2.400)--(2.920,1.590);
\draw[dashed] (-1.080,-0.810)--(0.000,0.000)--(4.000,0.000) (0.000,0.000)--(0.000,2.400);
\draw (4.000,0.000)--(4.000,2.400);
\draw (2.920,-0.810)--(4.000,0.000);
\node[below left] at (0.000,0.000){$O$};
\draw[->,gray] (-1.080,-0.810)--(-1.368,-1.026) node[below left]{$x$}; \draw[->,gray] (4.000,0.000)--(4.800,0.000) node[right]{$y$}; \draw[->,gray] (0.000,2.400)--(0.000,3.100) node[above]{$z$};
\draw (1.460,1.995)--(1.460,1.695); \shade[ball color=gray!30] (1.460,1.615) circle (.12); \node[right] at (1.460,1.615){đèn};
\filldraw[fill=gray!60] (2.920,1.590) circle (.08) node[right]{camera};

\end{tikzpicture}}

\pa{$14$.}{$3$.}{$7$.}{$11$.}
\dapso{C. Đèn ở chính giữa trần: $(2;2;3)$.}\end{bai}

\phanbt{Trắc nghiệm đúng/sai}
\begin{bai}Trong không gian $Oxyz$, cho $A(1;2;-1)$, $B(3;0;1)$.
\begin{dungsai}
  \y $\vt{AB}=(2;-2;2)$. \oDS
  \y Trung điểm của đoạn $AB$ là $I(2;1;0)$. \oDS
  \y Hình chiếu của $B$ lên trục $Oy$ là $(0;0;1)$. \oDS
  \y Điểm $A$ thuộc mặt phẳng $(Oxy)$. \oDS
\end{dungsai}
\dapso{a) Đúng; b) Đúng; c) Sai (là $(0;0;0)$); d) Sai.}\end{bai}
\begin{bai}Trong không gian $Oxyz$, cho điểm $M(2;-3;5)$.
\begin{dungsai}
  \y Hình chiếu của $M$ lên trục $Ox$ là $(2;0;0)$. \oDS
  \y Điểm đối xứng với $M$ qua mặt phẳng $(Oyz)$ là $(-2;-3;5)$. \oDS
  \y Điểm đối xứng với $M$ qua gốc toạ độ là $(-2;3;-5)$. \oDS
  \y Khoảng cách từ $M$ đến mặt phẳng $(Oxy)$ bằng $2$. \oDS
\end{dungsai}
\dapso{a) Đúng; b) Đúng; c) Đúng; d) Sai (bằng $|z|=5$).}\end{bai}

\begin{bai}Một phòng học có dạng hình hộp chữ nhật với chiều dài $6$ m, chiều rộng $5$ m, chiều cao $3$ m. Chọn hệ trục $Oxyz$ có gốc $O$ tại một góc sàn, các trục $Ox$, $Oy$, $Oz$ lần lượt theo chiều dài, chiều rộng, chiều cao của phòng. Một bóng đèn treo ở chính giữa trần nhà, một camera gắn ở góc trần đối diện với góc $O$ qua tâm phòng. Xét tính đúng sai của các mệnh đề sau:
\hinh{\begin{tikzpicture}[scale=0.8,font=\footnotesize,line join=round,line cap=round,>=Stealth]
\draw (-1.080,-0.810)--(2.920,-0.810)--(2.920,1.590)--(-1.080,1.590)--cycle;
\draw (-1.080,1.590)--(0.000,2.400)--(4.000,2.400)--(2.920,1.590);
\draw[dashed] (-1.080,-0.810)--(0.000,0.000)--(4.000,0.000) (0.000,0.000)--(0.000,2.400);
\draw (4.000,0.000)--(4.000,2.400);
\draw (2.920,-0.810)--(4.000,0.000);
\node[below left] at (0.000,0.000){$O$};
\draw[->,gray] (-1.080,-0.810)--(-1.368,-1.026) node[below left]{$x$}; \draw[->,gray] (4.000,0.000)--(4.800,0.000) node[right]{$y$}; \draw[->,gray] (0.000,2.400)--(0.000,3.100) node[above]{$z$};
\draw (1.460,1.995)--(1.460,1.695); \shade[ball color=gray!30] (1.460,1.615) circle (.12); \node[right] at (1.460,1.615){đèn};
\filldraw[fill=gray!60] (2.920,1.590) circle (.08) node[right]{camera};

\end{tikzpicture}}

\begin{dungsai}
  \y Bóng đèn cách camera $\dfrac{\sqrt{61}}{2}$ m. \oDS
  \y Góc giữa $OC$ và mặt sàn xấp xỉ $16^\circ$. \oDS
  \y Tổng các toạ độ của bóng đèn bằng $\dfrac{17}{2}$. \oDS
  \y Khoảng cách từ $O$ đến camera bằng $\dfrac{\sqrt{70}}{2}$ m. \oDS
\end{dungsai}
\dapso{$(3;\dfrac{5}{2};3)$, $(6;5;3)$; khoảng cách $=\dfrac{\sqrt{61}}{2}$. Camera tại $(6;5;3)$, khoảng cách $=\sqrt{6^2+5^2+3^2}=\sqrt{70}$. Đèn ở chính giữa trần: $(3;\dfrac{5}{2};3)$. $\overrightarrow{OC}=(6;5;3)$, $\sin\varphi=\dfrac{3}{\sqrt{70}}$, $\varphi\approx 21^\circ$. a) Đúng. b) Sai: giá trị đúng là $21$. c) Đúng. d) Sai: giá trị đúng là $\sqrt{70}$.}\end{bai}
\begin{bai}Một phòng học có dạng hình hộp chữ nhật với chiều dài $10$ m, chiều rộng $3$ m, chiều cao $3$ m. Chọn hệ trục $Oxyz$ có gốc $O$ tại một góc sàn, các trục $Ox$, $Oy$, $Oz$ lần lượt theo chiều dài, chiều rộng, chiều cao của phòng. Một bóng đèn treo ở chính giữa trần nhà, một camera gắn ở góc trần đối diện với góc $O$ qua tâm phòng. Xét tính đúng sai của các mệnh đề sau:
\hinh{\begin{tikzpicture}[scale=0.8,font=\footnotesize,line join=round,line cap=round,>=Stealth]
\draw (-1.080,-0.810)--(2.920,-0.810)--(2.920,1.590)--(-1.080,1.590)--cycle;
\draw (-1.080,1.590)--(0.000,2.400)--(4.000,2.400)--(2.920,1.590);
\draw[dashed] (-1.080,-0.810)--(0.000,0.000)--(4.000,0.000) (0.000,0.000)--(0.000,2.400);
\draw (4.000,0.000)--(4.000,2.400);
\draw (2.920,-0.810)--(4.000,0.000);
\node[below left] at (0.000,0.000){$O$};
\draw[->,gray] (-1.080,-0.810)--(-1.368,-1.026) node[below left]{$x$}; \draw[->,gray] (4.000,0.000)--(4.800,0.000) node[right]{$y$}; \draw[->,gray] (0.000,2.400)--(0.000,3.100) node[above]{$z$};
\draw (1.460,1.995)--(1.460,1.695); \shade[ball color=gray!30] (1.460,1.615) circle (.12); \node[right] at (1.460,1.615){đèn};
\filldraw[fill=gray!60] (2.920,1.590) circle (.08) node[right]{camera};

\end{tikzpicture}}

\begin{dungsai}
  \y Khoảng cách từ $O$ đến camera bằng $\sqrt{118}$ m. \oDS
  \y Tổng các toạ độ của bóng đèn bằng $\dfrac{21}{2}$. \oDS
  \y Bóng đèn cách camera $\dfrac{\sqrt{109}}{2}$ m. \oDS
  \y Góc giữa $OC$ và mặt sàn xấp xỉ $16^\circ$. \oDS
\end{dungsai}
\dapso{Đèn ở chính giữa trần: $(5;\dfrac{3}{2};3)$. $\overrightarrow{OC}=(10;3;3)$, $\sin\varphi=\dfrac{3}{\sqrt{118}}$, $\varphi\approx 16^\circ$. $(5;\dfrac{3}{2};3)$, $(10;3;3)$; khoảng cách $=\dfrac{\sqrt{109}}{2}$. Camera tại $(10;3;3)$, khoảng cách $=\sqrt{10^2+3^2+3^2}=\sqrt{118}$. a) Đúng. b) Sai: giá trị đúng là $\dfrac{19}{2}$. c) Đúng. d) Đúng.}\end{bai}

\phanbt{Trả lời ngắn}
\begin{bai}Một phòng học dạng hình hộp chữ nhật dài $8$ m, rộng $6$ m, cao $3$ m. Chọn hệ trục $Oxyz$ có gốc tại một góc sàn, các trục $Ox$, $Oy$, $Oz$ lần lượt theo chiều dài, chiều rộng, chiều cao (đơn vị mét). Đèn treo ở chính giữa trần nhà có toạ độ $(a;b;c)$. Tính $a+b+c$.\otraloi\dapso{$10$. Đèn ở $(4;3;3)$.}\end{bai}
\begin{bai}Gọi $M(a;b;c)$ là điểm đối xứng với $N(3;-1;2)$ qua trục $Oz$. Tính $a+b+c$.\otraloi\dapso{$0$. $M(-3;1;2)$.}\end{bai}
\begin{bai}Chọn hệ trục $Oxyz$ (đơn vị mét) có gốc tại trạm điều khiển, $Ox$ hướng đông, $Oy$ hướng bắc, $Oz$ hướng thẳng lên trời. Một drone từ trạm bay thẳng lên $50$ m, rồi bay về hướng đông $30$ m, rồi bay về hướng nam $40$ m. Toạ độ của drone là $(x;y;z)$. Tính $x+y+z$.\otraloi\dapso{$40$. Drone ở $(30;-40;50)$.}\end{bai}
\begin{bai}Cho hình hộp $ABCD.A'B'C'D'$ có $A(1;0;1)$, $B(2;1;2)$, $D(1;-1;1)$, $C'(4;5;-5)$. Biết $A'(a;b;c)$, tính $a+b+c$.\otraloi\dapso{$2$. $\vt{DC}=\vt{AB}\Rightarrow C(2;0;2)$; $\vt{AA'}=\vt{CC'}=(2;5;-7)\Rightarrow A'(3;5;-6)$.}\end{bai}
\begin{bai}Trong không gian $Oxyz$ (đơn vị: mét), mặt đất là mặt phẳng $(Oxy)$, hai flycam đang bay ở các vị trí $A(10;-10;7)$ và $B(-6;3;12)$. Tìm khoảng cách giữa hai flycam (đơn vị: $\text{ m}$) (làm tròn kết quả đến hàng phần mười).
\hinh{\begin{tikzpicture}[scale=0.85,font=\footnotesize,line join=round,line cap=round,>=Stealth]
\draw[->,gray] (0.000,0.000)--(-1.440,-1.080) node[below left]{$x$};
\draw[->,gray] (0.000,0.000)--(5.000,0.000) node[right]{$y$};
\draw[->,gray] (0.000,0.000)--(0.000,3.300) node[above]{$z$};
\node[below right,gray] at (0.000,0.000){$O$};
\fill[gray!10] (-0.500,0.000)--(-1.760,-0.945)--(3.240,-0.945)--(4.500,0.000)--cycle;
\draw[dashed] (0.300,1.525)--(0.300,-0.675); \fill (0.300,-0.675) circle (.04);
\filldraw[fill=gray!40] (0.300,1.525) ++(-.18,0) rectangle ++(.36,.1); \node[above] at (0.300,1.525){$A$};
\draw[dashed] (3.240,2.530)--(3.240,-0.270); \fill (3.240,-0.270) circle (.04);
\filldraw[fill=gray!40] (3.240,2.530) ++(-.18,0) rectangle ++(.36,.1); \node[above] at (3.240,2.530){$B$};
\draw[thick,dotted] (0.300,1.525)--(3.240,2.530);
\end{tikzpicture}}
\otraloi
\dapso{21,2. $AB=\sqrt{450}=15 \sqrt{2}$. Kết quả: $21{,}2$.}\end{bai}
\begin{bai}Trong không gian $Oxyz$ (đơn vị: mét), mặt đất là mặt phẳng $(Oxy)$, hai flycam đang bay ở các vị trí $A(-10;3;11)$ và $B(5;-3;8)$. Tìm khoảng cách từ flycam ở $B$ đến điểm $H$ trên mặt đất nằm ngay dưới $A$ (đơn vị: $\text{ m}$) (làm tròn kết quả đến hàng phần mười).
\hinh{\begin{tikzpicture}[scale=0.85,font=\footnotesize,line join=round,line cap=round,>=Stealth]
\draw[->,gray] (0.000,0.000)--(-1.440,-1.080) node[below left]{$x$};
\draw[->,gray] (0.000,0.000)--(5.000,0.000) node[right]{$y$};
\draw[->,gray] (0.000,0.000)--(0.000,3.300) node[above]{$z$};
\node[below right,gray] at (0.000,0.000){$O$};
\fill[gray!10] (-0.500,0.000)--(-1.760,-0.945)--(3.240,-0.945)--(4.500,0.000)--cycle;
\draw[dashed] (0.300,1.525)--(0.300,-0.675); \fill (0.300,-0.675) circle (.04);
\filldraw[fill=gray!40] (0.300,1.525) ++(-.18,0) rectangle ++(.36,.1); \node[above] at (0.300,1.525){$A$};
\draw[dashed] (3.240,2.530)--(3.240,-0.270); \fill (3.240,-0.270) circle (.04);
\filldraw[fill=gray!40] (3.240,2.530) ++(-.18,0) rectangle ++(.36,.1); \node[above] at (3.240,2.530){$B$};
\draw[thick,dotted] (0.300,1.525)--(3.240,2.530);
\end{tikzpicture}}
\otraloi
\dapso{18. $H(-10;3;0)$, $BH=5 \sqrt{13}$. Kết quả: $18$.}\end{bai}
\begin{bai}Một phòng học có dạng hình hộp chữ nhật với chiều dài $4$ m, chiều rộng $3$ m, chiều cao $4$ m. Chọn hệ trục $Oxyz$ có gốc $O$ tại một góc sàn, các trục $Ox$, $Oy$, $Oz$ lần lượt theo chiều dài, chiều rộng, chiều cao của phòng. Một bóng đèn treo ở chính giữa trần nhà, một camera gắn ở góc trần đối diện với góc $O$ qua tâm phòng. Tìm khoảng cách (m) từ bóng đèn đến camera.
\hinh{\begin{tikzpicture}[scale=0.8,font=\footnotesize,line join=round,line cap=round,>=Stealth]
\draw (-1.080,-0.810)--(2.920,-0.810)--(2.920,1.590)--(-1.080,1.590)--cycle;
\draw (-1.080,1.590)--(0.000,2.400)--(4.000,2.400)--(2.920,1.590);
\draw[dashed] (-1.080,-0.810)--(0.000,0.000)--(4.000,0.000) (0.000,0.000)--(0.000,2.400);
\draw (4.000,0.000)--(4.000,2.400);
\draw (2.920,-0.810)--(4.000,0.000);
\node[below left] at (0.000,0.000){$O$};
\draw[->,gray] (-1.080,-0.810)--(-1.368,-1.026) node[below left]{$x$}; \draw[->,gray] (4.000,0.000)--(4.800,0.000) node[right]{$y$}; \draw[->,gray] (0.000,2.400)--(0.000,3.100) node[above]{$z$};
\draw (1.460,1.995)--(1.460,1.695); \shade[ball color=gray!30] (1.460,1.615) circle (.12); \node[right] at (1.460,1.615){đèn};
\filldraw[fill=gray!60] (2.920,1.590) circle (.08) node[right]{camera};

\end{tikzpicture}}
\otraloi
\dapso{2,5. $(2;\dfrac{3}{2};4)$, $(4;3;4)$; khoảng cách $=\dfrac{5}{2}$. Kết quả: $2{,}5$.}\end{bai}
\phieutln

\phanbt{Tự luận}
\begin{bai}Trong không gian $Oxyz$, cho hình hộp chữ nhật $OABC.O'A'B'C'$ có $A(3;0;0)$, $C(0;5;0)$, $O'(0;0;2)$. Viết toạ độ các đỉnh của hình hộp; tìm toạ độ các vectơ $\vt{OB'}$, $\vt{A'C}$.
\dongke{7}
\dapso{$B(3;5;0)$, $A'(3;0;2)$, $B'(3;5;2)$, $C'(0;5;2)$; $\vt{OB'}=(3;5;2)$, $\vt{A'C}=(-3;5;-2)$.}\end{bai}
\begin{bai}Một chiếc bàn hình chữ nhật dài $1{,}2$ m, rộng $0{,}8$ m, cao $0{,}75$ m đặt sát góc phòng. Chọn hệ trục $Oxyz$ có gốc là góc phòng trên sàn, $Ox$, $Oy$ theo hai chân tường, $Oz$ thẳng đứng; hai cạnh mặt bàn song song với $Ox$, $Oy$ và mép bàn chạm hai bức tường. Tìm toạ độ bốn góc mặt bàn và toạ độ tâm mặt bàn.
\dongke{7}
\dapso{Các góc mặt bàn: $(0;0;0{,}75)$, $(1{,}2;0;0{,}75)$, $(1{,}2;0{,}8;0{,}75)$, $(0;0{,}8;0{,}75)$; tâm $(0{,}6;0{,}4;0{,}75)$.}\end{bai}
\begin{bai}Trong không gian $Oxyz$ (đơn vị: mét), mặt đất là mặt phẳng $(Oxy)$, hai flycam đang bay ở các vị trí $A(-1;-2;20)$ và $B(-8;-4;8)$. Tìm khoảng cách từ flycam ở $B$ đến điểm $H$ trên mặt đất nằm ngay dưới $A$ (đơn vị: $\text{ m}$) (làm tròn kết quả đến hàng phần mười).
\hinh{\begin{tikzpicture}[scale=0.85,font=\footnotesize,line join=round,line cap=round,>=Stealth]
\draw[->,gray] (0.000,0.000)--(-1.440,-1.080) node[below left]{$x$};
\draw[->,gray] (0.000,0.000)--(5.000,0.000) node[right]{$y$};
\draw[->,gray] (0.000,0.000)--(0.000,3.300) node[above]{$z$};
\node[below right,gray] at (0.000,0.000){$O$};
\fill[gray!10] (-0.500,0.000)--(-1.760,-0.945)--(3.240,-0.945)--(4.500,0.000)--cycle;
\draw[dashed] (0.300,1.525)--(0.300,-0.675); \fill (0.300,-0.675) circle (.04);
\filldraw[fill=gray!40] (0.300,1.525) ++(-.18,0) rectangle ++(.36,.1); \node[above] at (0.300,1.525){$A$};
\draw[dashed] (3.240,2.530)--(3.240,-0.270); \fill (3.240,-0.270) circle (.04);
\filldraw[fill=gray!40] (3.240,2.530) ++(-.18,0) rectangle ++(.36,.1); \node[above] at (3.240,2.530){$B$};
\draw[thick,dotted] (0.300,1.525)--(3.240,2.530);
\end{tikzpicture}}

\dongke{6}
\dapso{$H(-1;-2;0)$, $BH=3 \sqrt{13}$. Kết quả: $10{,}8$.}\end{bai}
\begin{bai}Một phòng học có dạng hình hộp chữ nhật với chiều dài $4$ m, chiều rộng $6$ m, chiều cao $3$ m. Chọn hệ trục $Oxyz$ có gốc $O$ tại một góc sàn, các trục $Ox$, $Oy$, $Oz$ lần lượt theo chiều dài, chiều rộng, chiều cao của phòng. Một bóng đèn treo ở chính giữa trần nhà, một camera gắn ở góc trần đối diện với góc $O$ qua tâm phòng. Tìm góc (độ) giữa đường thẳng nối góc $O$ với camera và mặt sàn, làm tròn đến hàng đơn vị,.
\hinh{\begin{tikzpicture}[scale=0.8,font=\footnotesize,line join=round,line cap=round,>=Stealth]
\draw (-1.080,-0.810)--(2.920,-0.810)--(2.920,1.590)--(-1.080,1.590)--cycle;
\draw (-1.080,1.590)--(0.000,2.400)--(4.000,2.400)--(2.920,1.590);
\draw[dashed] (-1.080,-0.810)--(0.000,0.000)--(4.000,0.000) (0.000,0.000)--(0.000,2.400);
\draw (4.000,0.000)--(4.000,2.400);
\draw (2.920,-0.810)--(4.000,0.000);
\node[below left] at (0.000,0.000){$O$};
\draw[->,gray] (-1.080,-0.810)--(-1.368,-1.026) node[below left]{$x$}; \draw[->,gray] (4.000,0.000)--(4.800,0.000) node[right]{$y$}; \draw[->,gray] (0.000,2.400)--(0.000,3.100) node[above]{$z$};
\draw (1.460,1.995)--(1.460,1.695); \shade[ball color=gray!30] (1.460,1.615) circle (.12); \node[right] at (1.460,1.615){đèn};
\filldraw[fill=gray!60] (2.920,1.590) circle (.08) node[right]{camera};

\end{tikzpicture}}

\dongke{6}
\dapso{$\overrightarrow{OC}=(4;6;3)$, $\sin\varphi=\dfrac{3}{\sqrt{61}}$, $\varphi\approx 23^\circ$. Kết quả: $23$.}\end{bai}

\phanbt{Bài tập tự luận rèn luyện}
\begin{bai}Một hộp quà có dạng hình hộp chữ nhật $OABC.O'A'B'C'$ được đặt vào hệ trục toạ độ $Oxyz$ sao cho đỉnh $O$ trùng gốc toạ độ, $A(3;0;0)$, $C(0;4;0)$ và $O'(0;0;2)$ (đơn vị: dm).
\hinh{\begin{tikzpicture}[scale=0.8,font=\footnotesize,line join=round,line cap=round,>=Stealth]
\draw[->,gray] (0.000,0.000)--(-1.296,-0.972) node[below left]{$x$};
\draw[->,gray] (0.000,0.000)--(5.000,0.000) node[right]{$y$};
\draw[->,gray] (0.000,0.000)--(0.000,2.900) node[above]{$z$};
\node[below right,gray] at (0.000,0.000){$O$};
\draw (-1.080,-0.810)--(2.920,-0.810)--(2.920,1.190)--(-1.080,1.190)--cycle;
\draw (-1.080,1.190)--(0.000,2.000)--(4.000,2.000)--(2.920,1.190);
\draw[dashed] (-1.080,-0.810)--(0.000,0.000)--(4.000,0.000) (0.000,0.000)--(0.000,2.000);
\draw (4.000,0.000)--(4.000,2.000);
\draw (2.920,-0.810)--(4.000,0.000);
\node[below] at (-1.080,-0.810){$A$};\node[below] at (4.000,0.000){$C$};\node[below] at (2.920,-0.810){$B$};\node[above] at (0.000,2.000){$O'$};\node[above] at (2.920,1.190){$B'$};
\end{tikzpicture}}
\begin{enumerate}[label=\alph*)]
\item Tìm toạ độ các đỉnh $B$, $A'$, $B'$, $C'$.
\item Tìm toạ độ của vectơ $\overrightarrow{OB'}$.
\item Tìm toạ độ hình chiếu vuông góc của $B'$ trên mặt phẳng $(Oxy)$ và trên trục $Oz$.
\item Tìm toạ độ trung điểm $M$ của đoạn $A'C'$.
\item Người ta gắn một chiếc nơ tại tâm $I$ của hộp quà (giao điểm các đường chéo). Tìm toạ độ của $I$.
\item Tìm toạ độ điểm đối xứng với $B'$ qua mặt phẳng $(Oxz)$.
\end{enumerate}
\dapso{a) $B(3;4;0)$, $A'(3;0;2)$, $B'(3;4;2)$, $C'(0;4;2)$. b) $(3;4;2)$. c) $(3;4;0)$ và $(0;0;2)$. d) $M\left(\dfrac32;2;2\right)$. e) $I\left(\dfrac32;2;1\right)$. f) $(3;-4;2)$.}\end{bai}

\begin{bai}Trong không gian $Oxyz$, cho hai vectơ $\vec a=2\vec i-3\vec j+\vec k$ và $\vec b=-\vec j+4\vec k$, cùng với điểm $M$ thoả mãn $\overrightarrow{OM}=\vec a$. Bài tập này giúp em làm quen với cách chuyển qua lại giữa dạng $x\vec i+y\vec j+z\vec k$ và dạng toạ độ.
\begin{enumerate}[label=\alph*)]
\item Viết toạ độ của hai vectơ $\vec a$ và $\vec b$.
\item Viết toạ độ của vectơ $\vec c=3\vec i-2\vec k$.
\item Tìm toạ độ của điểm $M$.
\item Tìm toạ độ hình chiếu vuông góc của $M$ trên mặt phẳng $(Oyz)$ và trên trục $Ox$.
\item Tính khoảng cách từ $M$ đến mặt phẳng $(Oxy)$ và đến trục $Oz$.
\item Tìm toạ độ điểm đối xứng với $M$ qua gốc toạ độ và qua trục $Ox$.
\end{enumerate}
\dapso{a) $\vec a=(2;-3;1)$, $\vec b=(0;-1;4)$. b) $(3;0;-2)$. c) $M(2;-3;1)$. d) $(0;-3;1)$; $(2;0;0)$. e) $1$; $\sqrt{13}$. f) $(-2;3;-1)$; $(2;3;-1)$.}\end{bai}

\begin{bai}Một trạm ra-đa đặt tại gốc $O$ theo dõi các máy bay trong vùng trời. Người ta chọn hệ trục $Oxyz$ có mặt phẳng $(Oxy)$ là mặt đất, trục $Ox$ hướng về phía đông, trục $Oy$ hướng về phía bắc, trục $Oz$ hướng thẳng đứng lên trời, đơn vị trên các trục là kilômét.
\hinh{\begin{tikzpicture}[scale=0.85,font=\footnotesize,line join=round,line cap=round,>=Stealth]
\draw[->,gray] (0.000,0.000)--(-1.440,-1.080) node[below left]{$x$};
\draw[->,gray] (0.000,0.000)--(5.000,0.000) node[right]{$y$};
\draw[->,gray] (0.000,0.000)--(0.000,3.300) node[above]{$z$};
\node[below right,gray] at (0.000,0.000){$O$};
\fill[gray!10] (-0.500,0.000)--(-1.760,-0.945)--(3.240,-0.945)--(4.500,0.000)--cycle;
\draw[dashed] (0.300,1.525)--(0.300,-0.675); \fill (0.300,-0.675) circle (.04);
\filldraw[fill=gray!40] (0.300,1.525) ++(-.18,0) rectangle ++(.36,.1); \node[above] at (0.300,1.525){$A$};
\draw[dashed] (3.240,2.530)--(3.240,-0.270); \fill (3.240,-0.270) circle (.04);
\filldraw[fill=gray!40] (3.240,2.530) ++(-.18,0) rectangle ++(.36,.1); \node[above] at (3.240,2.530){$B$};
\draw[thick,dotted] (0.300,1.525)--(3.240,2.530);
\end{tikzpicture}}
\begin{enumerate}[label=\alph*)]
\item Máy bay $A$ đang ở vị trí cách trạm $3$ km về phía đông, $5$ km về phía bắc và ở độ cao $2$ km. Viết toạ độ của $A$.
\item Máy bay $B$ ở vị trí cách trạm $4$ km về phía tây, $1$ km về phía nam, độ cao $1{,}5$ km. Viết toạ độ của $B$.
\item Tìm toạ độ của vectơ $\overrightarrow{AB}$.
\item Tìm toạ độ điểm $A'$ trên mặt đất nằm ngay phía dưới máy bay $A$.
\item Trong mỗi phút, máy bay $A$ dịch chuyển theo vectơ $\vec u=(0{,}5;-0{,}2;0{,}1)$. Tìm toạ độ của máy bay $A$ sau $1$ phút.
\item Nếu máy bay $A$ vẫn giữ nguyên cách dịch chuyển đó, sau $10$ phút nó ở vị trí nào và ở độ cao bao nhiêu?
\end{enumerate}
\dapso{a) $A(3;5;2)$. b) $B(-4;-1;1{,}5)$. c) $(-7;-6;-0{,}5)$. d) $A'(3;5;0)$. e) $(3{,}5;4{,}8;2{,}1)$. f) $(8;3;3)$, độ cao $3$ km.}\end{bai}

\begin{bai}Cho hình lập phương $ABCD.A'B'C'D'$ cạnh $2$. Chọn hệ trục $Oxyz$ có gốc trùng với $A$, các tia $Ox$, $Oy$, $Oz$ lần lượt đi qua $B$, $D$, $A'$.
\hinh{\begin{tikzpicture}[scale=0.9,font=\footnotesize,line join=round,line cap=round,>=Stealth]
\draw (-0.720,-0.540)--(1.280,-0.540)--(1.280,1.460)--(-0.720,1.460)--cycle;
\draw (-0.720,1.460)--(0.000,2.000)--(2.000,2.000)--(1.280,1.460);
\draw[dashed] (-0.720,-0.540)--(0.000,0.000)--(2.000,0.000) (0.000,0.000)--(0.000,2.000);
\draw (2.000,0.000)--(2.000,2.000);
\draw (1.280,-0.540)--(2.000,0.000);
\node[left] at (0.000,0.000){$A$};\node[below left] at (-0.720,-0.540){$B$};\node[below right] at (1.280,-0.540){$C$};\node[right] at (2.000,0.000){$D$};\node[left] at (0.000,2.000){$A'$};\node[left] at (-0.720,1.460){$B'$};\node[right] at (1.280,1.460){$C'$};\node[above right] at (2.000,2.000){$D'$};
\end{tikzpicture}}
\begin{enumerate}[label=\alph*)]
\item Viết toạ độ của tám đỉnh hình lập phương.
\item Tìm toạ độ tâm của hình lập phương.
\item Tìm toạ độ trung điểm của cạnh $B'C'$.
\item Tìm toạ độ của các vectơ $\overrightarrow{AC'}$ và $\overrightarrow{BD'}$.
\item Tìm toạ độ trọng tâm $G$ của tam giác $A'BD$.
\item Chứng minh $\overrightarrow{AG}=\dfrac13\overrightarrow{AC'}$, từ đó suy ra $G$ nằm trên đường chéo $AC'$.
\end{enumerate}
\dapso{a) $A(0;0;0)$, $B(2;0;0)$, $C(2;2;0)$, $D(0;2;0)$, $A'(0;0;2)$, $B'(2;0;2)$, $C'(2;2;2)$, $D'(0;2;2)$. b) $(1;1;1)$. c) $(2;1;2)$. d) $(2;2;2)$; $(-2;2;2)$. e) $G\left(\dfrac23;\dfrac23;\dfrac23\right)$. f) $\overrightarrow{AG}=\left(\dfrac23;\dfrac23;\dfrac23\right)=\dfrac13(2;2;2)$.}\end{bai}

\begin{bai}Lớp 12A học trong một phòng có dạng hình hộp chữ nhật dài $8$ m, rộng $6$ m, cao $3{,}5$ m. Để bố trí thiết bị, thầy chủ nhiệm chọn hệ trục $Oxyz$ có gốc $O$ là một góc sàn, các trục $Ox$, $Oy$, $Oz$ lần lượt chạy theo chiều dài, chiều rộng và chiều cao của phòng (đơn vị: mét).
\hinh{\begin{tikzpicture}[scale=0.8,font=\footnotesize,line join=round,line cap=round,>=Stealth]
\draw (-1.080,-0.810)--(2.920,-0.810)--(2.920,1.590)--(-1.080,1.590)--cycle;
\draw (-1.080,1.590)--(0.000,2.400)--(4.000,2.400)--(2.920,1.590);
\draw[dashed] (-1.080,-0.810)--(0.000,0.000)--(4.000,0.000) (0.000,0.000)--(0.000,2.400);
\draw (4.000,0.000)--(4.000,2.400);
\draw (2.920,-0.810)--(4.000,0.000);
\node[below left] at (0.000,0.000){$O$};
\draw[->,gray] (-1.080,-0.810)--(-1.368,-1.026) node[below left]{$x$}; \draw[->,gray] (4.000,0.000)--(4.800,0.000) node[right]{$y$}; \draw[->,gray] (0.000,2.400)--(0.000,3.100) node[above]{$z$};
\draw (1.460,1.995)--(1.460,1.695); \shade[ball color=gray!30] (1.460,1.615) circle (.12); \node[right] at (1.460,1.615){đèn};
\filldraw[fill=gray!60] (2.920,1.590) circle (.08) node[right]{camera};

\end{tikzpicture}}
\begin{enumerate}[label=\alph*)]
\item Viết toạ độ bốn góc trần nhà.
\item Máy chiếu được treo đúng ở chính giữa trần. Viết toạ độ của máy chiếu.
\item Màn chiếu treo trên bức tường nằm trong mặt phẳng $x=0$, tâm màn ở chính giữa bức tường theo chiều ngang và cách sàn $2$ m. Viết toạ độ tâm màn.
\item Tìm toạ độ của vectơ có điểm đầu là máy chiếu, điểm cuối là tâm màn.
\item Một camera an ninh gắn ở góc trần đối diện với góc $O$. Viết toạ độ camera và toạ độ hình chiếu vuông góc của nó trên sàn.
\item Mắt của bạn Hoa khi ngồi học ở vị trí $H(5;2;1{,}2)$. Cho biết mắt bạn ấy cách sàn bao nhiêu và cách bức tường $x=0$ bao nhiêu mét.
\end{enumerate}
\dapso{a) $(0;0;3{,}5)$, $(8;0;3{,}5)$, $(8;6;3{,}5)$, $(0;6;3{,}5)$. b) $(4;3;3{,}5)$. c) $(0;3;2)$. d) $(-4;0;-1{,}5)$. e) $(8;6;3{,}5)$ và $(8;6;0)$. f) $1{,}2$ m và $5$ m.}\end{bai}
